\section{概率论基础}\label{ux6982ux7387ux8bbaux57faux7840}

\subsection{Set Theory（集合论）}\label{set-theoryux96c6ux5408ux8bba}

\subsubsection{From random experiment to probability model（从随机试验到概率模型）}\label{from-random-experiment-to-probability-modelux4eceux968fux673aux8bd5ux9a8cux5230ux6982ux7387ux6a21ux578b}

A random experiment has well-defined possible outcomes, but the realized outcome is uncertain before the experiment.

\[
\text{random experiment} \to \text{sample space } S \to \text{events } A \to \text{probabilities } P(A)
\]

（随机试验 → 样本空间 → 事件 → 概率）

\subsubsection{Sample space and events（样本空间与事件）}\label{sample-space-and-eventsux6837ux672cux7a7aux95f4ux4e0eux4e8bux4ef6}

\begin{itemize}
\tightlist
\item
  \textbf{Sample space} \(S\)（样本空间）: the set of all possible basic outcomes; finite（有限）, countably infinite（可数无限）, or uncountable（不可数）.
\item
  \textbf{Countable}（可数）: can be placed in one-to-one correspondence with a subset of the integers.
\item
  \textbf{Event}（事件）: any subset \(A\subseteq S\); \(A\) occurs iff the realized outcome \(\in A\).
\item
  \(\varnothing\): the impossible event（不可能事件）.
\item
  \(A\subseteq B\iff (x\in A\Rightarrow x\in B)\); \(A=B\iff A\subseteq B\) and \(B\subseteq A\). Prove set equality by double inclusion（双向包含）.
\end{itemize}

\subsubsection{Set operations and algebra of sets（集合运算与代数性质）}\label{set-operations-and-algebra-of-setsux96c6ux5408ux8fd0ux7b97ux4e0eux4ee3ux6570ux6027ux8d28}

\[
A\cup B=\{x:x\in A\text{ or }x\in B\},\qquad
A\cap B=\{x:x\in A\text{ and }x\in B\},
\]
\[
A^{c}=\{x\in S:x\notin A\}.
\]

\begin{itemize}
\tightlist
\item
  Commutative, associative, distributive laws all hold;
\item
  De Morgan's laws: \((A\cup B)^{c}=A^{c}\cap B^{c}\), \((A\cap B)^{c}=A^{c}\cup B^{c}\);
\item
  Generalized De Morgan for countable unions/intersections:
  \[
  \left(\bigcup_i A_i\right)^{c}=\bigcap_i A_i^{c},\qquad
  \left(\bigcap_i A_i\right)^{c}=\bigcup_i A_i^{c}.
  \]
\item
  \(\bigcup_{i=1}^{\infty}A_i\): belongs to at least one \(A_i\); \(\bigcap_{i=1}^{\infty}A_i\): belongs to every \(A_i\).
\end{itemize}

\subsubsection{Partitions（分割）}\label{partitionsux5206ux5272}

\begin{itemize}
\tightlist
\item
  \textbf{Mutually exclusive} (pairwise disjoint)（两两互斥）: \(A_i\cap A_j=\varnothing\) for \(i\ne j\);
\item
  \textbf{Collectively exhaustive}（穷举）: \(\bigcup_i A_i=S\);
\item
  \textbf{Partition}（分割）: pairwise disjoint and collectively exhaustive. Central for the law of total probability and Bayes' rule.
\end{itemize}

\subsection{Basics of Probability Theory（概率论基础）}\label{basics-of-probability-theoryux6982ux7387ux8bbaux57faux7840}

\subsubsection{\texorpdfstring{\(\sigma\)-algebra（\(\sigma\)-代数）}{\textbackslash sigma-algebra（\textbackslash sigma-代数）}}\label{sigma-algebrasigma-ux4ee3ux6570}

A collection \(\mathcal{B}\) of subsets of \(S\) is a \textbf{\(\sigma\)-algebra} if:

\begin{enumerate}
\def\labelenumi{\arabic{enumi}.}
\tightlist
\item
  \(\varnothing\in\mathcal{B}\);
\item
  \(A\in\mathcal{B}\Rightarrow A^{c}\in\mathcal{B}\);
\item
  \(A_1,A_2,\dots\in\mathcal{B}\Rightarrow\bigcup_{i=1}^{\infty}A_i\in\mathcal{B}\).
\end{enumerate}

By De Morgan's laws, a \(\sigma\)-algebra is also closed under countable intersections. The power set \(2^{S}\) and the trivial family \(\{\varnothing,S\}\) are both \(\sigma\)-algebras.

\subsubsection{Probability function and axioms（概率函数与公理）}\label{probability-function-and-axiomsux6982ux7387ux51fdux6570ux4e0eux516cux7406}

A complete probability model is \((S,\mathcal{B},P)\), with \(P:\mathcal{B}\to[0,1]\):

\begin{enumerate}
\def\labelenumi{\arabic{enumi}.}
\tightlist
\item
  Nonnegativity（非负性）: \(P(A)\ge0\);
\item
  Normalization（规范性）: \(P(S)=1\);
\item
  Countable additivity（可数可加性）: for pairwise disjoint \(A_1,A_2,\dots\),
  \[
  P\left(\bigcup_{i=1}^{\infty}A_i\right)=\sum_{i=1}^{\infty}P(A_i).
  \]
\end{enumerate}

\textbf{Key point}: the axioms regulate how probabilities behave; they do not determine which probability model fits a given experiment（公理只约束概率运算，不决定具体模型；那是建模假设）.

\subsubsection{Consequences of the axioms（公理推论）}\label{consequences-of-the-axiomsux516cux7406ux63a8ux8bba}

\[
P(\varnothing)=0,\qquad 0\le P(A)\le1,\qquad P(A^{c})=1-P(A),
\]
\[
A\subseteq B\Rightarrow P(A)\le P(B)\quad\text{(monotonicity)}.
\]

Two-event identities（两事件公式）:

\[
P(B\cap A^{c})=P(B)-P(A\cap B),
\]
\[
P(A\cup B)=P(A)+P(B)-P(A\cap B).
\]

\subsubsection{Partition identity and inequalities（分割公式与不等式）}\label{partition-identity-and-inequalitiesux5206ux5272ux516cux5f0fux4e0eux4e0dux7b49ux5f0f}

\begin{itemize}
\tightlist
\item
  If \(\{C_i\}\) is a partition of \(S\), then（全概率结构）:
  \[
  P(A)=\sum_i P(A\cap C_i).
  \]
\item
  \textbf{Boole's inequality}:
  \[
  P\left(\bigcup_i A_i\right)\le\sum_i P(A_i).
  \]
\item
  \textbf{Bonferroni lower bound} (via Boole on complements):
  \[
  P\left(\bigcap_{i=1}^{n}A_i\right)\ge\sum_{i=1}^{n}P(A_i)-(n-1).
  \]
  Two-event case: \(P(A\cap B)\ge P(A)+P(B)-1\).
\end{itemize}

\subsubsection{Constructing a probability on a countable space（可数样本空间上构造概率）}\label{constructing-a-probability-on-a-countable-spaceux53efux6570ux6837ux672cux7a7aux95f4ux4e0aux6784ux9020ux6982ux7387}

Let \(S=\{s_1,s_2,\dots\}\) be finite or countable. If \(p_i\ge0\) and \(\sum_i p_i=1\), then

\[
P(A)=\sum_{i:\,s_i\in A}p_i
\]

defines a probability function: assign probability mass \(p_i\) to each outcome and sum over the event.

\subsubsection{Counting（计数）}\label{countingux8ba1ux6570}

\begin{itemize}
\tightlist
\item
  Multiplication principle（分步计数原理）: \(k\) successive tasks with \(n_1,\dots,n_k\) choices give \(n_1n_2\cdots n_k\) ways.
\item
  Factorial（阶乘）: \(n!=n(n-1)\cdots1\), with \(0!=1\).
\end{itemize}

Sampling \(r\) objects from \(n\):

\begin{longtable}[]{@{}
  >{\raggedright\arraybackslash}p{(\linewidth - 4\tabcolsep) * \real{0.3333}}
  >{\raggedright\arraybackslash}p{(\linewidth - 4\tabcolsep) * \real{0.3333}}
  >{\raggedright\arraybackslash}p{(\linewidth - 4\tabcolsep) * \real{0.3333}}@{}}
\toprule\noalign{}
\begin{minipage}[b]{\linewidth}\raggedright
Sampling
\end{minipage} & \begin{minipage}[b]{\linewidth}\raggedright
Ordered
\end{minipage} & \begin{minipage}[b]{\linewidth}\raggedright
Unordered
\end{minipage} \\
\midrule\noalign{}
\endhead
\bottomrule\noalign{}
\endlastfoot
Without replacement（不放回） & \(\dfrac{n!}{(n-r)!}\) & \(\dbinom{n}{r}=\dfrac{n!}{r!(n-r)!}\) \\
With replacement（放回） & \(n^{r}\) & \(\dbinom{n+r-1}{r}\) \\
\end{longtable}

\begin{itemize}
\tightlist
\item
  Stars and bars（隔板法）: arrange \(r\) stars and \(n-1\) bars, then choose \(r\) positions for the stars.
\item
  Multinomial coefficient（多重集排列）: with \(k\) positions and values repeated \(k_1,\dots,k_m\) times,
  \[
  \frac{k!}{k_1!\,k_2!\cdots k_m!}.
  \]
\end{itemize}

\subsubsection{Equally likely outcomes（等可能模型）}\label{equally-likely-outcomesux7b49ux53efux80fdux6a21ux578b}

If the finite sample space has \(N\) equally likely outcomes,

\[
P(A)=\frac{|A|}{|S|}.
\]

\textbf{Warning}: valid only when the elementary outcomes are truly equally likely. Probabilities must follow the actual sampling mechanism; unordered outcomes are not automatically equally likely（不能为了方便把无序结果当作等可能）.

\subsection{Conditional Probability and Independence（条件概率与独立性）}\label{conditional-probability-and-independenceux6761ux4ef6ux6982ux7387ux4e0eux72ecux7acbux6027}

\subsubsection{Conditional Probability（条件概率）}\label{conditional-probabilityux6761ux4ef6ux6982ux7387}

\textbf{Definition}

对事件 \(A,B\)，若 \(P(B)>0\)，则在 \(B\) 已发生的条件下，\(A\) 的条件概率为

\[
P(A\mid B)=\frac{P(A\cap B)}{P(B)}.
\]

获知 \(B\) 已发生后，新的有效样本空间被限制到 \(B\)；\(P(A\mid B)\) 衡量 \(B\) 中同时属于 \(A\) 的概率质量占比。

\subsubsection{Multiplication rule（乘法公式）}\label{multiplication-ruleux4e58ux6cd5ux516cux5f0f}

由定义直接得到

\[
P(A\cap B)=P(A\mid B)P(B)=P(B\mid A)P(A).
\]

对 \(A_1,\ldots,A_n\) 反复使用乘法公式，可得 \textbf{chain rule}：

\[
P\left(\bigcap_{i=1}^{n}A_i\right)
=P(A_1)\prod_{i=2}^{n}
P\left(A_i\,\middle|\,\bigcap_{j=1}^{i-1}A_j\right),
\]

前提是式中的条件概率均有定义。

\subsubsection{Example：连续抽到四张 A}\label{exampleux8fdeux7eedux62bdux5230ux56dbux5f20-a}

一副 52 张的标准扑克牌不放回抽取，前四张全是 A 的概率为

\[
P(\text{first four cards are aces})
=\frac{4}{52}\cdot\frac{3}{51}\cdot\frac{2}{50}\cdot\frac{1}{49}
=\frac{1}{\binom{52}{4}}.
\]

每次抽牌都会改变下一次的条件样本空间，因此各因子的分子、分母依次减少。

\subsubsection{Example：不放回抽球}\label{exampleux4e0dux653eux56deux62bdux7403}

盒中有 \(r\) 个红球和 \(b\) 个蓝球。连续不放回抽取两个球，令

\[
A=\{\text{first ball is red}\},
\qquad
B=\{\text{second ball is blue}\}.
\]

则

\[
P(A)=\frac{r}{r+b},
\qquad
P(B\mid A)=\frac{b}{r+b-1},
\]

所以

\[
P(A\cap B)=\frac{r}{r+b}\cdot\frac{b}{r+b-1}.
\]

\subsubsection{Law of Total Probability（全概率公式）}\label{law-of-total-probabilityux5168ux6982ux7387ux516cux5f0f}

若 \(A_1,A_2,\ldots\) 构成样本空间 \(S\) 的一个 partition，即两两互斥且并集为 \(S\)，则对任意事件 \(B\)，

\[
P(B)=\sum_i P(B\cap A_i)
=\sum_i P(B\mid A_i)P(A_i).
\]

它把事件 \(B\) 按照所有可能来源 \(A_i\) 拆开，再将各路径的概率相加。

\subsubsection{Example：三个组别}\label{exampleux4e09ux4e2aux7ec4ux522b}

若 \(P(B_i)=1/3\)，且 \(P(A\mid B_i)=i/6\)，\(i=1,2,3\)，则

\[
P(A)
=\sum_{i=1}^{3}P(A\mid B_i)P(B_i)
=\frac{1}{3}\left(\frac{1}{6}+\frac{2}{6}+\frac{3}{6}\right)
=\frac{1}{3}.
\]

\subsubsection{Bayes' Rule（贝叶斯公式）}\label{bayes-ruleux8d1dux53f6ux65afux516cux5f0f}

若 \(A_1,A_2,\ldots\) 构成 partition 且 \(P(B)>0\)，则

\[
P(A_i\mid B)
=\frac{P(B\mid A_i)P(A_i)}
{\sum_j P(B\mid A_j)P(A_j)}.
\]

\textbf{Interpretation}

Bayes' rule 把条件方向从 \(P(B\mid A_i)\) 反转为 \(P(A_i\mid B)\)：

\[
\text{posterior}
\propto
\text{likelihood}\times\text{prior}.
\]

分母是由全概率公式得到的 \(P(B)\)，负责把各个未标准化权重归一化。

\subsubsection{Example：汽车保险风险分类}\label{exampleux6c7dux8f66ux4fddux9669ux98ceux9669ux5206ux7c7b}

客户被分为高风险 \(H\)、中风险 \(M\) 和低风险 \(L\)：

\[
P(H)=0.25,\qquad P(M)=0.25,\qquad P(L)=0.50,
\]

一年内至少发生一次事故的条件概率为

\[
P(A\mid H)=0.25,\qquad
P(A\mid M)=0.16,\qquad
P(A\mid L)=0.10.
\]

已知某客户发生事故，则其属于高风险客户的概率为

\[
\begin{aligned}
P(H\mid A)
&=\frac{P(A\mid H)P(H)}
{P(A\mid H)P(H)+P(A\mid M)P(M)+P(A\mid L)P(L)}\\
&=\frac{0.25\times0.25}
{0.25\times0.25+0.16\times0.25+0.10\times0.50}\\
&\approx0.410.
\end{aligned}
\]

观察到事故后，高风险类型的概率从 prior \(0.25\) 更新到 posterior \(0.410\)。

\subsubsection{Example：色盲问题}\label{exampleux8272ux76f2ux95eeux9898}

设总体中男性人数与女性人数之比为 \(r:1\)，则

\[
P(M)=\frac{r}{1+r},
\qquad
P(W)=\frac{1}{1+r}.
\]

若色盲事件为 \(D\)，且

\[
P(D\mid M)=p,
\qquad
P(D\mid W)=p^2,
\]

则随机选出的色盲者为男性的概率为

\[
P(M\mid D)
=\frac{p\frac{r}{1+r}}
{p\frac{r}{1+r}+p^2\frac{1}{1+r}}
=\frac{r}{r+p}.
\]

\subsubsection{Monty Hall Problem（三门问题）}\label{monty-hall-problemux4e09ux95e8ux95eeux9898}

奖品等概率地藏在帘子 \(A,B,C\) 后。先选择 \(A\)，主持人知道奖品位置，并在另外两扇帘子中打开一扇空帘；假设他打开了 \(B\)。记

\[
B^*=\{\text{host opens curtain }B\text{ and it is empty}\}.
\]

主持人的行为规则意味着

\[
P(B^*\mid A)=\frac{1}{2},
\qquad
P(B^*\mid B)=0,
\qquad
P(B^*\mid C)=1.
\]

由 Bayes' rule，

\[
P(A\mid B^*)
=\frac{\frac12\cdot\frac13}
{\frac12\cdot\frac13+0\cdot\frac13+1\cdot\frac13}
=\frac13,
\]

\[
P(C\mid B^*)=\frac23.
\]

\textbf{Conclusion}

换到 \(C\) 的获奖概率为 \(2/3\)，高于坚持选择 \(A\) 的 \(1/3\)。这个结论依赖主持人``必定打开未选择的空帘，并在有两扇空帘可选时等概率选择''的行为规则。

\subsubsection{Independence（独立性）}\label{independenceux72ecux7acbux6027}

\textbf{Definition}

事件 \(A\) 与 \(B\) 独立，当且仅当

\[
P(A\cap B)=P(A)P(B).
\]

若 \(P(B)>0\)，则等价地有

\[
P(A\mid B)=P(A),
\]

即得知 \(B\) 发生不会改变 \(A\) 的概率。

\subsubsection{Independence vs.~disjointness（独立与互斥）}\label{independence-vs.-disjointnessux72ecux7acbux4e0eux4e92ux65a5}

二者含义不同：

\begin{itemize}
\tightlist
\item
  \textbf{Disjointness}：\(A\cap B=\varnothing\)，两个事件不能同时发生；
\item
  \textbf{Independence}：一个事件是否发生不提供关于另一个事件的信息。
\end{itemize}

若 \(A,B\) 互斥且 \(P(A)>0,P(B)>0\)，则

\[
P(A\cap B)=0<P(A)P(B),
\]

所以它们不可能独立。只有当至少一个事件概率为 \(0\) 时，互斥与独立才可能同时成立。

\subsubsection{Example：四次掷骰至少出现一次 6}\label{exampleux56dbux6b21ux63b7ux9ab0ux81f3ux5c11ux51faux73b0ux4e00ux6b21-6}

独立掷一枚公平骰子四次。使用 complement rule，

\[
\begin{aligned}
P(\text{at least one six})
&=1-P(\text{no six})\\
&=1-\left(\frac56\right)^4\\
&=\frac{671}{1296}\approx0.518.
\end{aligned}
\]

更一般地，若 \(A_1,\ldots,A_n\) 相互独立，则其补事件也相互独立，从而

\[
P\left(\bigcup_{i=1}^{n}A_i\right)
=1-P\left(\bigcap_{i=1}^{n}A_i^c\right)
=1-\prod_{i=1}^{n}\bigl(1-P(A_i)\bigr).
\]

若 \(P(A_i)=p\)，则

\[
P\left(\bigcup_{i=1}^{n}A_i\right)=1-(1-p)^n.
\]

\subsubsection{Review}\label{review}

\begin{enumerate}
\def\labelenumi{\arabic{enumi}.}
\tightlist
\item
  \textbf{Conditional probability} 是在已知信息下重新校准概率模型。
\item
  \textbf{Multiplication rule} 计算一条联合路径；\textbf{total probability} 汇总所有互斥路径。
\item
  \textbf{Bayes' rule} 用观测证据把 prior 更新为 posterior。
\item
  \textbf{Independence} 表示信息不改变概率，不能与``互斥''混淆。
\end{enumerate}

\subsubsection{Independence and complements（独立性与补事件）}\label{independence-and-complementsux72ecux7acbux6027ux4e0eux8865ux4e8bux4ef6}

\textbf{Theorem}

若 \(A\) 与 \(B\) 独立，则以下三对事件也分别独立：

\[
(A,B^c),\qquad (A^c,B),\qquad (A^c,B^c).
\]

例如，利用 \(A=(A\cap B)\cup(A\cap B^c)\) 这一互斥分解，

\[
\begin{aligned}
P(A\cap B^c)
&=P(A)-P(A\cap B)\\
&=P(A)-P(A)P(B)\\
&=P(A)\bigl[1-P(B)\bigr]\\
&=P(A)P(B^c).
\end{aligned}
\]

其余两种情形可用相同方法证明。因此，独立事件的任意补集组合仍保持独立。

\subsubsection{Pairwise vs.~mutual independence（两两独立与相互独立）}\label{pairwise-vs.-mutual-independenceux4e24ux4e24ux72ecux7acbux4e0eux76f8ux4e92ux72ecux7acb}

当事件数量不少于三个时，\textbf{pairwise independence} 不足以推出 \textbf{mutual independence}。

三个事件两两独立只要求

\[
P(A\cap B)=P(A)P(B),\qquad
P(A\cap C)=P(A)P(C),\qquad
P(B\cap C)=P(B)P(C),
\]

但不保证

\[
P(A\cap B\cap C)=P(A)P(B)P(C).
\]

\textbf{Definition}

事件 \(A_1,\ldots,A_n\) 相互独立，是指对任意子集
\(\{i_1,\ldots,i_k\}\subseteq\{1,\ldots,n\}\)，都有

\[
P\left(\bigcap_{j=1}^{k}A_{i_j}\right)
=\prod_{j=1}^{k}P(A_{i_j}).
\]

\begin{quote}
Mutual independence 要求所有规模至少为 \(2\) 的子事件组都满足乘法关系；只检查每一对事件并不充分。
\end{quote}

\subsubsection{Example：两两独立但不相互独立}\label{exampleux4e24ux4e24ux72ecux7acbux4f46ux4e0dux76f8ux4e92ux72ecux7acb}

设

\[
S=\{a,b,c,d\},
\]

四个结果等可能，并定义

\[
A_1=\{a,b\},\qquad
A_2=\{b,c\},\qquad
A_3=\{a,c\}.
\]

每个事件的概率均为 \(1/2\)，任意两个事件的交集均只含一个结果，因此

\[
P(A_i\cap A_j)=\frac14
=\frac12\cdot\frac12
=P(A_i)P(A_j),\qquad i\ne j.
\]

所以三个事件两两独立。然而

\[
A_1\cap A_2\cap A_3=\varnothing,
\]

从而

\[
P(A_1\cap A_2\cap A_3)=0
\ne\frac18
=P(A_1)P(A_2)P(A_3).
\]

故它们并不相互独立。

\subsubsection{Quiz check（课堂练习核对）}\label{quiz-checkux8bfeux5802ux7ec3ux4e60ux6838ux5bf9}

\begin{enumerate}
\def\labelenumi{\arabic{enumi}.}
\tightlist
\item
  盒中有 \(4\) 个红球、\(6\) 个蓝球，不放回抽取两个球：
\end{enumerate}

\[
P(\text{first red, second blue})
=\frac4{10}\cdot\frac6{9}
=\frac4{15}.
\]

\begin{enumerate}
\def\labelenumi{\arabic{enumi}.}
\setcounter{enumi}{1}
\tightlist
\item
  保险分类例中，
\end{enumerate}

\[
P(A)=0.25(0.25)+0.16(0.25)+0.10(0.50)=0.1525,
\]

\[
P(H\mid A)=\frac{0.25(0.25)}{0.1525}
=\frac{25}{61}\approx0.410.
\]

\begin{enumerate}
\def\labelenumi{\arabic{enumi}.}
\setcounter{enumi}{2}
\tightlist
\item
  上述 \(A_1,A_2,A_3\) 的每个两两交集概率为 \(1/4\)，但三者交集概率为 \(0\)，因此两两独立而不相互独立。
\end{enumerate}

\subsection{Random Variables（随机变量）}\label{random-variablesux968fux673aux53d8ux91cf}

\textbf{Definition}

随机变量 \(X\) 是从样本空间到实数集的函数：

\[
X:S\longrightarrow\mathbb{R}.
\]

\(X\) 的所有可能取值构成其值域，通常记为 \(\mathcal X\)。随机变量并不是``随机变化的数''，而是先把每个原始结果 \(s\in S\) 映射为数值 \(X(s)\)；随机性来自试验结果 \(s\) 的不确定性。

\[
s\in S
\quad\xrightarrow{\ X\ }\quad
X(s)\in\mathbb R.
\]

\subsubsection{Examples（随机变量示例）}\label{examplesux968fux673aux53d8ux91cfux793aux4f8b}

\begin{longtable}[]{@{}ll@{}}
\toprule\noalign{}
Random experiment & Random variable \\
\midrule\noalign{}
\endhead
\bottomrule\noalign{}
\endlastfoot
掷两枚骰子 & \(X=\) 两个点数之和 \\
抛硬币 25 次 & \(X=\) 正面出现次数 \\
测量作物反应 & \(X=\) 单位面积产量 \\
抽样调查 50 位选民 & \(X=\) 支持某提案的人数 \\
\end{longtable}

随机变量可以把复杂的原始结果压缩为当前分析所关心的 numerical summary。

\subsubsection{Induced probabilities（随机变量诱导的概率）}\label{induced-probabilitiesux968fux673aux53d8ux91cfux8bf1ux5bfcux7684ux6982ux7387}

给定 \(X:S\to\mathcal X\)，样本空间 \(S\) 上的概率会诱导出 \(X\) 取值上的概率。对 \(x\in\mathcal X\)，

\[
P(X=x)
=P\bigl(\{s\in S:X(s)=x\}\bigr).
\]

更一般地，对 \(A\subseteq\mathcal X\)，

\[
P(X\in A)
=P\bigl(\{s\in S:X(s)\in A\}\bigr).
\]

右侧集合是 \(A\) 在映射 \(X\) 下的原像。因此，对随机变量事件求概率，本质上仍然是在样本空间 \(S\) 上求概率。

\textbf{Notation}

\begin{itemize}
\tightlist
\item
  大写字母 \(X,Y,Z\) 表示随机变量；
\item
  小写字母 \(x,y,z\) 表示随机变量的可能取值或实现值。
\end{itemize}

\subsubsection{Example：三次抛硬币}\label{exampleux4e09ux6b21ux629bux786cux5e01}

公平硬币独立抛三次，令 \(X\) 为正面次数。样本空间为

\[
S=\{HHH,HHT,HTH,THH,HTT,THT,TTH,TTT\}.
\]

\(X\) 的概率分布为

\begin{longtable}[]{@{}rrrrr@{}}
\toprule\noalign{}
\(x\) & \(0\) & \(1\) & \(2\) & \(3\) \\
\midrule\noalign{}
\endhead
\bottomrule\noalign{}
\endlastfoot
\(P(X=x)\) & \(1/8\) & \(3/8\) & \(3/8\) & \(1/8\) \\
\end{longtable}

因此

\[
P(0\le X\le1)=P(X=0)+P(X=1)=\frac48,
\]

而

\[
P(0\le X<1)=P(X=0)=\frac18.
\]

对离散型随机变量，把 \(\le\) 改成 \(<\) 可能改变事件本身，不能忽略端点。

\subsection{Distribution Functions（分布函数）}\label{distribution-functionsux5206ux5e03ux51fdux6570}

\subsubsection{Cumulative distribution function（累积分布函数）}\label{cumulative-distribution-functionux7d2fux79efux5206ux5e03ux51fdux6570}

\textbf{Definition}

随机变量 \(X\) 的 cumulative distribution function（CDF）定义为

\[
F_X(x)=P(X\le x),\qquad x\in\mathbb R.
\]

CDF 在所有实数 \(x\) 上都有定义，而不只是在 \(X\) 的值域 \(\mathcal X\) 上定义。若 \(a<b\)，则

\[
P(a<X\le b)=F_X(b)-F_X(a).
\]

\subsubsection{Example：三次抛硬币的 CDF}\label{exampleux4e09ux6b21ux629bux786cux5e01ux7684-cdf}

对正面次数 \(X\)，

\[
F_X(x)=
\begin{cases}
0, & x<0,\\
\dfrac18, & 0\le x<1,\\
\dfrac12, & 1\le x<2,\\
\dfrac78, & 2\le x<3,\\
1, & x\ge3.
\end{cases}
\]

这是一个阶梯函数。每个取值点 \(x_0\) 处的跳跃大小等于该点的概率质量：

\[
F_X(x_0)-F_X(x_0^-)=P(X=x_0),
\]

其中

\[
F_X(x_0^-)=\lim_{x\uparrow x_0}F_X(x).
\]

CDF 在跳跃点取跳跃后的函数值，因此是 \textbf{right-continuous}（右连续）的。

\subsubsection{Example：命中点到圆心的距离}\label{exampleux547dux4e2dux70b9ux5230ux5706ux5fc3ux7684ux8dddux79bb}

射击命中半径为 \(R\) 的圆形靶，且命中点在靶面上按面积均匀分布。令 \(X\) 为命中点到圆心的距离。

当 \(0\le x\le R\) 时，事件 \(\{X\le x\}\) 对应半径为 \(x\) 的内圆，因此

\[
F_X(x)
=\frac{\pi x^2}{\pi R^2}
=\frac{x^2}{R^2}.
\]

完整的 CDF 为

\[
F_X(x)=
\begin{cases}
0, & x<0,\\
\dfrac{x^2}{R^2}, & 0\le x\le R,\\
1, & x>R.
\end{cases}
\]

\textbf{Key point}

``在圆面上按面积均匀''不等于``距离 \(X\) 在 \([0,R]\) 上均匀''。半径越大的环带面积越大，因此 \(F_X(x)\) 是二次函数而不是 \(x/R\)。

\subsubsection{Review}\label{review-1}

\begin{enumerate}
\def\labelenumi{\arabic{enumi}.}
\tightlist
\item
  两个独立事件取补后仍独立；三个及以上事件必须区分 pairwise independence 与 mutual independence。
\item
  随机变量是函数 \(X:S\to\mathbb R\)，它把原始结果转换为可分析的数值。
\item
  随机变量事件的概率由其在样本空间中的原像决定。
\item
  CDF \(F_X(x)=P(X\le x)\) 在整个实数轴上定义，并且右连续。
\item
  离散分布中，CDF 的跳跃大小恰好等于对应点的概率质量。
\end{enumerate}
